\section*{Chapitre \ref{ch:rc:portfolio-credit} --- Crédit de portefeuille}

\begin{solution}{exo:rc:portfolio-credit:1}
$125\times4.08~\% = 5.10$ défauts. Non~: par linéarité de l'espérance, le nombre attendu de
défauts ne dépend que des lois marginales~; la corrélation modifie la dispersion autour de cette valeur.
\end{solution}

\begin{solution}{exo:rc:portfolio-credit:2}
La tranche equity est longue en corrélation~: une corrélation plus élevée accroît la probabilité d'avoir peu
de défauts, c'est là qu'elle survit. Les tranches seniors sont courtes en corrélation~: seuls
les scénarios corrélés les atteignent. La mezzanine se situe entre les deux~: à faible corrélation, une
corrélation plus forte y fait entrer des pertes, à forte corrélation les pertes la dépassent d'un bond ou n'ont pas lieu,
si bien que sa valeur n'est pas monotone.
\end{solution}

\begin{solution}{exo:rc:portfolio-credit:3}
Chaque défaut coûte $0.6/125 = 0.48~\%$ du portefeuille~: la tranche equity est épuisée dès le septième
défaut ($3/0.48 = 6.25$), et les pertes atteignent 22~\% au quarante-sixième ($22/0.48 = 45.8$).
\end{solution}

\begin{solution}{exo:rc:portfolio-credit:4}
La \omterm{def:rc:portfolio-credit:dc}{corrélation de défaut} est la corrélation des indicateurs d'événements rares. Deux variables latentes
de corrélation 25~\% évoluent souvent ensemble, mais toutes deux doivent passer sous un seuil de 1.74
écart-type sous leur moyenne pour un défaut conjoint~; la probabilité conjointe s'élève nettement au-dessus de
$p^2$, mais reste faible par rapport à $p$, si bien que la corrélation des indicateurs est faible.
\end{solution}

\begin{solution}{exo:rc:portfolio-credit:5}
$50\times3.28 = \text{EUR}~164$ millions de protection achetée sur l'indice.
\end{solution}

\begin{solution}{exo:rc:portfolio-credit:6}
Les valeurs des tranches de base sont monotones en corrélation, si bien qu'une corrélation de base existe toujours et
est unique et interpole entre les points de détachement~; les \omterm{def:rc:portfolio-credit:compound}{corrélations composées} peuvent valoir deux ou
aucune. Mais valoriser la mezzanine comme la différence de deux tranches de base à des
corrélations différentes n'est pas un modèle~: des corrélations de base interpolées peuvent donner des pertes attendues
négatives pour les tranches fines, et les portefeuilles sur mesure et les maturités exigent une règle de correspondance.
\end{solution}

\begin{solution}{exo:rc:portfolio-credit:7}
\texttt{tail\_table()}~: 0.13~\% sous la \omterm{def:rc:portfolio-credit:gauss}{copule gaussienne} et 1.13~\% sous la \omterm{def:rc:portfolio-credit:t}{copule de Student} avec
quatre degrés de liberté, soit un rapport de 8.6.
\end{solution}

\begin{solution}{exo:rc:portfolio-credit:8}
La perte attendue n'est pas le risque~: la valeur d'une \omterm{def:rc:portfolio-credit:ss}{tranche super senior} varie avec la
probabilité de pertes extrêmes et corrélées, qui croît en crise précisément quand le détenteur
ne peut pas vendre~; son notionnel élevé transforme de faibles variations de prix en lourdes pertes~; et une notation
fondée sur la queue de distribution d'un modèle hérite de son erreur dans cette zone. Elle exige une couverture ou du capital pour
sa valeur de marché, comme l'a montré 2007--08.
\end{solution}

\begin{solution}{pb:rc:portfolio-credit:1}
\textbf{1.} 17.3 et 7.32 fois l'indice~; 23.56 millions d'euros sur la tranche 3--6~\%.
\textbf{2.} Upfront de 3.721 millions d'euros~; en running, 500 points de base sur la tranche equity (500\,000~euros
par an), moins 195 points de base sur la mezzanine (459\,502~euros), soit environ 40\,000~euros par an tant qu'aucun
défaut ne survient.
\textbf{3.} Neutre à un léger élargissement parallèle~; longue au risque de défaut de tout nom individuel (la
tranche equity absorbe les premières pertes)~; longue en corrélation.
\textbf{4.} Elle gagne quand la corrélation monte (la valeur de la tranche equity monte davantage que celle de la
mezzanine).
\textbf{5.} Variations idiosyncratiques, variations de corrélation, convexité des grands mouvements et
évolution des deltas dans le temps.
\textbf{6.} Upfront de la tranche equity 42.79~\%~; mezzanine 236.5 points de base.
\textbf{7.} Tranche equity $-\text{EUR}~557\,577$, mezzanine $+\text{EUR}~424\,041$, total
$-\text{EUR}~133\,536$.
\textbf{8.} Le delta porte sur tous les noms évoluant ensemble~; l'élargissement d'un seul nom déplace la
tranche equity bien davantage que la mezzanine, par rapport à leurs deltas parallèles.
\textbf{9.} 0.48~\% du portefeuille, 16~\% de la tranche equity~: 1.6~millions d'euros.
\textbf{10.} Non~: l'indice bouge de la part d'un seul nom, $1/125$, tandis que la tranche equity supporte une grande
part du risque de ce nom~; la protection sur ce nom est la couverture.
\textbf{11.} Tranche equity $-\text{EUR}~269\,618$, mezzanine $-\text{EUR}~594\,873$ (son spread tombe à 138
points de base), total $-\text{EUR}~864\,491$.
\textbf{12.} La perte attendue 3--6~\% vaut $2\,\mathrm{EL}_{0,6}-\mathrm{EL}_{0,3}$~: à $\rho_6$ inchangé, un $\rho_3$
plus faible fait monter $\mathrm{EL}_{0,3}$ et abaisse celle de la mezzanine.
\textbf{13.} Tranche equity $-\text{EUR}~828\,391$, mezzanine $-\text{EUR}~171\,869$, total $-\text{EUR}~1\,000\,261$.
\textbf{14.} La protection sur la mezzanine, achetée pour gagner en cas d'élargissement~; son spread a baissé avec
la corrélation, si bien qu'elle a aussi perdu.
\textbf{15.} Les vendeurs de protection equity la rachètent (élargissement de la tranche equity) et revendent leur
protection sur la mezzanine (resserrement de celle-ci), ce qui fait jouer les deux jambes contre chaque position semblable.
\textbf{16.} Dans le modèle, la corrélation d'actifs de la \omterm{def:rc:portfolio-credit:gauss}{copule gaussienne}~; sur le marché, le
prix de la tranche equity par rapport au reste de la structure, fixé par l'offre et la
demande.
\textbf{17.} Limites sur la sensibilité à la corrélation et sur le saut au défaut par nom, tests
de résistance sur les variations conjointes idiosyncratiques et de corrélation, et la liquidité des couvertures.
\textbf{18.} La \omterm{def:rc:portfolio-credit:gauss}{copule gaussienne}, calibrée sur des marchés calmes, avait des queues fines, et
les corrélations des portefeuilles hypothécaires étaient bien plus élevées que supposé~; les modèles mesuraient
des prix de marché, non le risque du facteur.
\textbf{19.} Cas nommé~: \emph{le correlation trade de mai 2005}~: la jambe equity perd
828\,391~euros et la couverture de la mezzanine 171\,869~euros, soit au total 1\,000\,261~euros sur
10~millions d'euros de tranche equity.
\textbf{20.} Sur la corrélation et sur la dispersion des noms~: le delta n'élimine que les
variations parallèles de spread.
\end{solution}

\begin{solution}{iq:rc:portfolio-credit:1}
$X_i = \sqrt\rho Z+\sqrt{1-\rho}\varepsilon_i$, défaut d'ici $t$ si $X_i < N^{-1}(p_i(t))$~;
conditionnellement à $Z = z$, $p_i(t\mid z) = N((N^{-1}(p_i(t))-\sqrt\rho z)/\sqrt{1-\rho})$, et
les noms sont indépendants.
\iqlookfor{la variable latente, le seuil et l'indépendance conditionnelle.}
\end{solution}

\begin{solution}{iq:rc:portfolio-credit:2}
Conditionner au facteur~; les noms sont alors indépendants, et la loi des pertes
s'obtient par la récursion qui ajoute un nom à la fois (sur une grille de pertes pour des notionnels inégaux)~;
intégrer sur le facteur par quadrature de Gauss--Hermite. Coût~: nombre de noms fois grille de pertes fois
nombre de nœuds de quadrature, par date.
\iqlookfor{indépendance conditionnelle, récursion et quadrature.}
\end{solution}

\begin{solution}{iq:rc:portfolio-credit:3}
Composée~: une corrélation par tranche~; les tranches mezzanines peuvent en avoir deux ou aucune. De base~: la
corrélation de la tranche de base $[0,d]$, toujours unique~; les tranches sont évaluées comme des différences de
tranches de base à des corrélations différentes, ce qui peut créer un arbitrage (pertes espérées négatives) et
exige une mise en correspondance pour les portefeuilles sur mesure.
\iqlookfor{non-unicité contre incohérence.}
\end{solution}

\begin{solution}{iq:rc:portfolio-credit:4}
Par la valeur de marché, et non par la perte espérée~: sensibilités aux spreads, à la corrélation et au facteur~;
tests de résistance de scénarios systémiques~; modèles de queue de distribution au-delà du gaussien~; risque de
contrepartie de toute protection achetée (monolines, AIG)~; limites de liquidité et de concentration.
\iqlookfor{risque en valeur de marché, tests de résistance et contreparties de couverture.}
\end{solution}

\begin{solution}{iq:rc:portfolio-credit:5}
La limite de la probabilité qu'une variable soit extrême sachant que l'autre l'est. La
\omterm{def:rc:portfolio-credit:gauss}{copule gaussienne} n'en a aucune (pour une corrélation inférieure à un)~; la \omterm{def:rc:portfolio-credit:t}{copule de Student} et la copule de Clayton
en ont une. Les tranches seniors ne subissent des pertes que lorsque de nombreux noms font défaut ensemble, ce que
régit précisément la dépendance de queue.
\iqlookfor{la définition et le lien avec les tranches seniors.}
\end{solution}

\begin{solution}{iq:rc:portfolio-credit:6}
Il était neutre face aux mouvements parallèles, mais long du risque de nom unique et long de la corrélation. Les
explosions idiosyncratiques (GM et Ford) ont élargi la tranche equity bien plus que la mezzanine~; la corrélation
implicite dans les cours des actions a baissé, ce qui a fait baisser davantage la tranche equity et le spread de la
mezzanine, de sorte que la protection de la mezzanine a elle aussi perdu de la valeur~; les dénouements ont amplifié les deux mouvements.
\iqlookfor{risque idiosyncratique, risque de corrélation et effet de foule.}
\end{solution}
